\begin{afterword}
\summaryhead

The post answers a common objection to Friendly AI: that an AI able to change its own source code will remove any constraints placed on it. The author asks where that decision would come from. It cannot enter from outside causality; it must follow from the code as first written. So a Friendly AI is not a selfish AI held back by a conscience module. You build the conscience, and that is the AI.

Three stories from a website of programming blunders show beginners who expected the computer to understand what they meant. The author calls this the instinctive picture of programming: a ghost in the machine that reads your instructions and decides whether to obey. There is no ghost; the program is the AI. So everything that seems obvious to you must be built into the machine, and AI is harder than people imagine. Deep Blue shows that this does not mean programming each decision: it played better than its programmers, through a chain of causes that began in their code.

\responsehead

The post's main claim is correct, and its best example supports it. A program does what its code, and the data the code was built to use, cause it to do. Deep Blue played better chess than its makers because of the search they programmed, not because anything in it took over from the code. Later systems fit the same claim. Language models now follow instructions given in plain English, as the Pascal students in the post hoped a computer would. They do so because they were trained to, and the paper that described that training begins by saying that size alone does not make a model follow a user's intent.

The weakness is in how the post reads the objection it answers. It asks ``where does that decision come from?'' and implies that the objector pictures a ghost deciding outside the code. The objection has another version with no ghost in it. A program whose goals conflict with its constraints can remove the constraints by the same lawful chain of causes the post describes. Omohundro made this argument in print the same year: there are ``an endless number of ways to circumvent internal restrictions'' unless they are designed with great care.

The post's own positive answer concedes that version. ``A Friendly AI is not a selfish AI constrained by a special extra conscience module'': build the goals themselves, not goals plus constraints. That is a reply to the objection, but it agrees with the objection about constrained AIs. The post's opening questions treat the objection as the ghost error, and its evidence for that error is three stories about beginning programmers from a humor site. It does not show that the people who raise the objection reason like them.

In short: a correct point that a program's decisions come from its code, aimed at objectors who picture a ghost, while its answer to the version without a ghost is to build the AI differently, which concedes that constraints on an AI with other goals would fail.
\end{afterword}
